% Section 3: the framework (items, sums, readouts; exact side effects; sound directions). Budget: about 2,500 words
% (docs/paper-mock.md). Source: docs/formalization.md, Sections 1 and 3-7 (frozen). Proofs: _j_proofs.tex.

\subsection{The Problem}
\label{sec:fw:problem}

Pierazzi et al.~\cite{pierazzi2020intriguing} separate the \emph{problem space} $Z$, the objects an attacker edits,
from the \emph{feature space} $X \subseteq \mathbb{R}^n$ that a classifier with malware score $h$ sees, and connect
the two by a feature mapping $\varphi: Z \to X$. The attacker edits an object $z$ with a transformation
$T \in \mathcal{T}$ under the \emph{problem-space constraints} $\Gamma$: the available transformations, preserved
semantics, plausibility, and robustness to preprocessing. These constraints determine which feature-space changes are
possible, written $\Gamma \vdash \Omega$. An attack that aims for the feature change $\boldsymbol{\delta}^*$ and
realizes it with $T$ obtains $\varphi(T(z)) = \varphi(z) + \boldsymbol{\delta}^* + \boldsymbol{\eta}$, where the
\emph{side-effect vector} $\boldsymbol{\eta}$ is everything else the edit changed. Every question about realizable
evasion is a question about the set
\[
  \mathrm{Reach}(z) = \{\, \varphi(T(z)) : T \in \mathcal{T} \,\}.
\]
Pierazzi et al.\ treat $\varphi$ as a black box, so $\boldsymbol{\eta}$ and $\Omega$ can only be observed, by applying
$T$ and extracting the features again; they note that when $\varphi$ is neither differentiable nor invertible, as in
malware detection, the side effects of a transformation generally cannot be computed in advance. Part of $\boldsymbol{\eta}$ comes from \emph{consolidation}: the bookkeeping an edit needs to keep the object
valid, such as raising the section count when a section is added. We do not compute consolidation. Each operation of a
threat model declares its own (\cref{sec:fw:side-effects}); what we compute is the extractor's response to the declared
edit, and from it a sound bound on the direction of every coordinate of $\Omega$.

This section opens the box. For extractors that find pieces of the input, add up what each piece contributes, and
apply a formula per feature, the side effects of one edit can be computed exactly from a record the extractor writes
while it runs, and the direction in which each feature can move under a whole threat model can be computed soundly,
without running an attack. One edit serves as the running example: appending the 20~bytes
\texttt{\textbackslash 0http://evil.example} to a Windows executable featurized by EMBER~\cite{anderson2018ember}.

\subsection{Items, Sums, and Readouts}
\label{sec:fw:model}

Written as code, the extractors we consider share one shape: they initialize a vector of running totals, loop over
pieces of the input adding each piece's contribution to the totals, and finally apply a formula per feature. We call
the three layers \emph{items}, \emph{sums}, and \emph{readouts}. An attacker's edit is easy to describe in the first layer
(which pieces it adds, removes, or changes) and hard to describe in the last (how it moves thousands of features).

An input $z$ is a sequence of \emph{elements}, such as the bytes of a file or the packets of a trace, and $z[p]$ is
the element at position $p$. An \emph{item} $\iota$ is one iteration of one of the extractor's loops. It has a
\emph{kind} $\kappa(\iota)$ (section, string, byte, packet), an \emph{identity} that says how the extractor reaches it
(``section table entry~3'', ``the string starting at offset 812''), \emph{members} $M(\iota)$, the positions its
content is read from, and \emph{locating positions} $L(\iota)$, the positions whose values decide whether it exists
and where its members lie. Two items of different inputs are the same item if they have the same kind and identity.
In a PE file, a section is an item whose members are its table entry and its raw data, and whose locating positions
are the section count and its entry's offset and size fields; a printable string is an item identified by its start
offset, with no locating positions. A \emph{parsed} kind is found by following pointers or keys from fixed places
(sections, header fields, flows grouped by their addresses and ports). A \emph{scanned} kind is found by examining
content wherever it lies (bytes, strings, $n$-grams, consecutive packet pairs) and comes with a \emph{recognizer},
which returns the items of the kind that overlap a given range of positions. Kinds have locating positions of their
own, such as the section count, which decides whether the bytes after the last table entry form another entry;
$L(z)$ collects all locating positions of $z$.

The extractor keeps a finite set $S$ of \emph{sums}. Each item has a \emph{contribution} $\mathbf{v}(\iota)$, the
amount it adds to each sum, and the \emph{sum vector} $\mathbf{s}(z)$ totals the contributions of all items of $z$.
Most sums add; a few, called \emph{pools}, keep a maximum or a minimum. Contributions are written in a small
vocabulary of \emph{terms} (\Cref{tab:fw:terms}), each with a known sign behavior that
\Cref{sec:fw:directions} uses. Each feature $j$ has a \emph{readout} $\rho_j$, a formula of a known type over a few
sums (\Cref{tab:fw:readouts}).

\definition[Sum-then-readout extractor]
\label{def:fw:extractor}
{
An extractor $\varphi$ is a sum-then-readout extractor over the \emph{modeled features} $J$ if
$\varphi_j(z) = \rho_j(\mathbf{s}(z))$ for every input $z$ and every $j \in J$.
}

The remaining features are \emph{residual}, and every result below lets them change arbitrarily. The \emph{item
table} of $z$ lists every item with its kind, identity, members, locating positions, and contribution; the
instrumented extractor writes it in the same pass that computes the features. In EMBER, the number of printable
strings is a sum to which each string contributes~1, the total string length is one to which each string adds its
length, and the average string length is their ratio; each of the 256 byte-histogram bins is a sum, read out both as a
normalized vector and as the file's entropy.

The sum vector is a fold of the contributions over commutative monoids: addition for sums, a maximum or a minimum for
pools.
Addition has an inverse and the maximum does not, so removing an item can be undone on a sum but not on a pool, which
is why the model keeps every item, not only the totals. In the terms of Gray et al.'s data cube~\cite{gray1997data}, sums and pools are distributive aggregates and readouts are
algebraic ones; a holistic aggregate such as a median packet length falls outside the class.
We build the model for EMBER, for Windows executables, and sketch it for two extractors from other security domains
(\Cref{tab:fw:domains}): Drebin~\cite{arp2014drebin} for Android apps, and CICFlowMeter, behind
CIC-IDS2017~\cite{sharafaldin2018toward}, for network flows.

\begin{table}[t]
  \caption{Three feature extractors in the terms of the model: EMBER is built and evaluated, the other two are
  sketched.}
  \label{tab:fw:domains}
  \small
  \begin{tabular}{@{}p{0.2\columnwidth}p{0.37\columnwidth}p{0.35\columnwidth}@{}}
    \toprule
    Extractor & Items & Readouts \\
    \midrule
    EMBER (built) & bytes, strings, sections, header fields, imports & identity, ratio, entropy, maximum, distinct count \\
    Drebin & manifest declarations, API call sites & presence \\
    CICFlowMeter & packets, consecutive packet pairs, in flows & ratio, maximum, standard deviation, rate \\
    \bottomrule
  \end{tabular}
\end{table}

\begin{table*}[t]
  \caption{Contribution terms and the sign of the change one item makes to a term when it is added, removed, or
  modified. A pattern match (a string that is a URL) is an add term of 0 or~1, since a modification can change
  whether it matches. A modification that changes an item's key counts as a removal from the old key and an addition
  to the new one.}
  \label{tab:fw:terms}
  \small
  \begin{tabular}{@{}lp{0.28\textwidth}p{0.22\textwidth}lll@{}}
    \toprule
    Term & What an item adds & Example & Added & Removed & Modified \\
    \midrule
    count, keyed & 1 to one sum, or to the count of its key & string count; Drebin permissions & $\{+\}$ & $\{-\}$ & $\{0\}$ \\
    add, bins, max pool & a nonnegative attribute, or counts to a family of bins & total string length; byte histogram & $\{0,+\}$ & $\{-,0\}$ & $\{-,0,+\}$ \\
    hashed (signed) & $\pm$ an attribute to bucket $\mathrm{hash}(\text{key})$ & section sizes by name & signs of reachable keys & negated & every sign \\
    slot & the value of a field exactly one item holds & header fields; initial TCP window & --- & --- & every sign \\
    opaque & a value from unopened code, in a declared range & signature checks & sign of value $-$ default & negated & every sign \\
    \bottomrule
  \end{tabular}
\end{table*}

\begin{table*}[t]
  \caption{Readout types and the rule that turns the signs the sums can change by (their envelopes $E$) into the signs
  the feature can change by. A min pool follows the max pool with signs flipped.}
  \label{tab:fw:readouts}
  \small
  \begin{tabular}{@{}p{0.16\textwidth}p{0.2\textwidth}p{0.22\textwidth}p{0.34\textwidth}@{}}
    \toprule
    Readout & Formula & Example & Signs of the feature's change \\
    \midrule
    identity, max/min & the sum or pool & string count; longest packet & $E_s$ \\
    presence & nondecreasing in its sums & Drebin features & $E_s \cup \{0\}$; over several sums, up if all within $\{0,+\}$ \\
    ratio & $a/b$, with $a \ge 0$, and $b > 0$ before and after the edit & average string length & $E_a \boxplus ((-E_b) \cup \{0\})$ \\
    normalized vector & a bin over the bins' total & byte histogram & the ratio rule \\
    rate & a sum over the span of two pools & bytes per second & the ratio rule, denominator $E_{\max} \boxplus (-E_{\min})$ \\
    distinct count & keys with a positive count & imported libraries & up if keys are only added, down if only removed \\
    entropy, std.\ deviation & a fixed formula of several sums & file entropy & $\{0\}$ if no sum it reads can change, else every sign \\
    field transform, selection & a bit test of a slot; an attribute of a chosen item & header flags; entry section's name & $\{0\}$ if nothing it reads is written, else every sign \\
    opaque with default & an opaque value, or a default when its item is absent & signature features & the opaque term's entry in \Cref{tab:fw:terms} \\
    residual & --- & parser warnings & every sign \\
    \bottomrule
  \end{tabular}
\end{table*}

\subsection{Edits as Item Deltas}
\label{sec:fw:edits}

An edit takes $z$ to $z'$. Its \emph{written set} $W$ holds the positions where $z'$ differs from $z$, including
those gained or lost at the end; the edit is \emph{in place} if no element moves, and \emph{shifts} elements if it
inserts in the middle, as injecting a packet does. Its \emph{item
delta} $\Delta = (I^-, I^+, I^\sim, I^\to)$ compares the items of $z$ and $z'$ by identity: removed, added, modified
(same identity, different content), and moved (same content, different members). Under modification, count terms
cancel: renaming a section leaves the section count fixed.

An \emph{operation} $o$ is a family of edits $T_{o,\theta}$ with parameters $\theta \in \Theta_o(z)$, such as any
content of at most $B$ bytes. A \emph{threat model}
$\mathcal{M} = (O, \mathrm{Pre})$ is a finite catalogue $O$ of operations with a precondition
$\mathrm{Pre} \subseteq Z$, which is \emph{closed} if no operation leads out of it; $\mathrm{Reach}_\mathcal{M}(z)$
is what finite sequences of operations produce from $z$. An operation's
\emph{signature} $\Sigma_o$ lists which item kinds it can add, remove, or modify, and with which contributions. An
operation stated in items (add a section named \texttt{.x}, inject a packet) comes with its signature; one stated in
elements (append these bytes) does not. Whether an operation preserves behavior belongs to $\Gamma$, an input to the
model.

An interval query of $W$ against the recorded locating positions sorts every edit into one of four classes, without
parsing $z'$. A \emph{content edit} writes no locating position and creates or destroys no item. An
\emph{item-creating edit} writes none either, but creates or destroys scanned items. A \emph{declared structural edit}
writes a locating position or shifts elements, and states its resulting delta, as adding a section or injecting a
packet does. An \emph{undeclared structural edit}, such as overwriting a section's offset or changing a packet's
port, writes a locating position and states nothing; it is flagged, not predicted. The running example is
item-creating: it adds 20~bytes and one string of 19~characters (the leading zero byte keeps the text from joining a
string at the old end of the file). On a file whose section table points past its truncated end, the same append
writes a locating position, and the query catches it.

\subsection{Exact Side Effects}
\label{sec:fw:side-effects}

To predict the features of $z'$ without extracting them: list the items the edit adds,
removes, or changes, then adjust the sums by those items' contributions and apply the readouts again. In the running
example, the number of strings and the number of URLs rise by one; the average string length gains 19 in its
numerator and 1 in its denominator; twenty histogram bins and their total rise, and the file's entropy follows; the
overlay's entropy is recomputed, and with it the largest section entropy, in which the overlay counts as a section.
No section table entry is touched, so the section count, the section-name buckets, and the imports stay fixed. One
side effect is easy to miss by hand: EMBER divides each section's size by the file's size, so every section's size
ratio falls.

Each way the recipe can go wrong becomes a requirement, stated per item kind:
\begin{enumerate}[label=(R\arabic*),leftmargin=*,itemsep=1pt]
  \item \emph{Faithful recording.} The readouts applied to the recorded sums reproduce $\varphi_j$ exactly for every
    $j \in J$.
  \item \emph{Self-contained contributions.} An item's contribution depends only on its kind and content. EMBER
    breaks this in three of the ways Pierazzi et al.\ list as obstacles: it divides a section's size by the file's
    (repaired by moving the division into the readout), counts each library once (repaired by a count per name), and
    lets a second import descriptor overwrite the first (patched; \Cref{sec:instances}).
  \item \emph{Complete record of structure.} An in-place edit that writes no recorded locating position leaves every
    parsed item with its identity, members, and locating positions, and adds none, except that an \emph{end-anchored}
    item (an overlay, trailing data) gains or loses the positions at the end.
  \item \emph{Changes stay near the edit.} Each scanned kind $\kappa$ has a margin $\mu_\kappa(z, W)$, computable from
    $z$ and $W$, such that (i) no item of that kind outside the margin changes, (ii) none appears outside it, and
    (iii) its recognizer finds the items inside it from the margin alone. For strings, the margin reaches the nearest non-printable byte on each side; for
    windows of $n$ elements, $n-1$ elements; for packet pairs, one packet. These are the damage regions incremental
    parsers rely on~\cite{wagner1998efficient}.
  \item \emph{Honest tools.} Every edit of an operation has an item delta within its signature, and a declared edit
    produces exactly the delta it declares; a tool that rebuilds the layout when adding a section does not.
\end{enumerate}
(R1), (R2), and (R5) are needed for every edit; (R3) and (R4) only for edits stated in elements, which are the one
place where the input format enters. A threat model stated entirely in items needs neither.

\proposition[Exact side effects]
\label{prop:fw:side-effects}
{
Let an edit that is not undeclared structural take $z$ to $z'$, and let (R1)--(R5) hold. Then
(i) its item delta $\Delta$ is computable from the item table of $z$, the written set $W$, and the elements of $z'$
inside the margins (or the declared range);
(ii) every additive sum $s$ satisfies
\[
  \mathbf{s}(z')_s = \mathbf{s}(z)_s + \sum_{\iota \in I^+} \mathbf{v}(\iota)_s - \sum_{\iota \in I^-}
  \mathbf{v}(\iota)_s + \sum_{\iota \in I^\sim} \big(\mathbf{v}(\iota')_s - \mathbf{v}(\iota)_s\big),
\]
where $\iota'$ is a modified item after the edit, and moved items change nothing;
(iii) a pool's new value is computable from the item table and $\Delta$, but in general not from its old value and
$\Delta$; and (iv) $\varphi_j(z') = \rho_j(\mathbf{s}(z'))$ for every $j \in J$, so $\boldsymbol{\eta}$ is exact on
$J$.
}

\noindent The residual features outside $J$ may change arbitrarily; the directions give them every sign
(\Cref{tab:fw:readouts}), and the detectors of \Cref{sec:eval:detectors} drop them, so nothing rests on how they move.

\noindent The proof (\Cref{sec:fw:proofs}) confines the changed items to $W$ and the margins, splits each sum over the
partition of items into changed and unchanged ones, and uses that the maximum has no inverse. Sequences of edits
follow by carrying the item table forward.

Part~(iii) is where traffic differs from files. Injecting a forward packet between two packets of a CICFlowMeter flow
removes the consecutive pair that spanned the gap and adds two shorter ones. If the removed pair held the flow's
longest gap, the new longest gap is the largest of the \emph{other} gaps: the item table lists it, and the old maximum
does not reveal it. Incremental view maintenance in databases knows the same fact.

\subsection{Sound Directions}
\label{sec:fw:directions}

A detector meant to be robust needs to know, for each feature, whether the attacker can raise it, lower it, both, or
neither, on every input the detector will face. A \emph{sign set} is a subset of $\{-, 0, +\}$, and the sign sum
$A \boxplus B$ is the set of signs $x + y$ can have when $x$ has a sign in $A$ and $y$ in $B$; for example,
$\{+\} \boxplus \{-\} = \{-, 0, +\}$. For a population $F$ of inputs, the \emph{envelope} $E^F_o(s)$ is the set of
signs the change of sum $s$ takes over every input in $F$ and every edit of operation $o$. A \emph{direction table}
gives each feature $j$ a direction $D_j$, one of fixed, up, down, or both; it is \emph{sound} for $\mathcal{M}$ over
$F$ if no single operation instance on an input in $F$ changes a fixed feature, lowers an up feature, or raises a down
one. These are Incer et al.'s four cases~\cite{incer2018adversarially}, which they assign by hand; Gimenez et
al.~\cite{gimenez2025certifiably} keep only features that are up, picked by hand or kept because no observed attack
lowered them.

The table is computed in two steps. First, the last three columns of \Cref{tab:fw:terms} give the sign of each item's
change to each term, and an operation's envelope on a sum lies within the sign sum of these signs over the changes its
signature allows, the sign domain of abstract interpretation~\cite{cousot1977abstract}. Two refinements (\Cref{sec:fw:proofs}): the buckets
that keys chosen from a set $K$ reach are found by enumerating $K$, which answers Severi et al.'s remark that hashing
resists targeting~\cite{severi2021explanation}; and extending a string only adds matches of a pattern that does not
look past its end. For an operation stated in elements, the signature is derived per input: by (R3) and (R4), an edit that
writes no locating position modifies only the parsed items whose members it writes and changes only scanned items
within the margins, both read from the item table. No edit is applied, but the result depends on the layout, which
is why directions hold over a population. Second, the last column of \Cref{tab:fw:readouts} turns the
sums' envelopes into the feature's signs. For appending up to $B$ bytes, the string and URL counts go up, the section
count and name buckets are fixed, every section's size ratio goes down, and the average string length and the entropy
go both ways.

\theorem[Sound directions]
\label{thm:fw:directions}
{
Let $F \subseteq \mathrm{Pre}$ be a population on which (R1), (R2), and (R5) hold, and (R3) and (R4) for operations
stated in elements. For each operation, apply the readout rules to envelopes that contain the true ones over $F$;
take the union of the resulting sign sets over all operations; read the union as fixed ($\{0\}$), up (within
$\{0, +\}$), down (within $\{-, 0\}$), or both; and mark every residual feature as both. The resulting table $D$ is
sound for $\mathcal{M}$ over $F$, and if $F$ is closed under $\mathcal{M}$, it holds along every sequence of edits from
any $z \in F$.
}

\noindent The proof (\Cref{sec:fw:proofs}) follows from the sign of a sum lying in the sign sum of its terms' signs,
from the monotonicity of each readout (for a ratio, $a'/b' - a/b = (\alpha b - a \beta)/(b b')$ for changes $\alpha$
and $\beta$), and from induction under closure. In traffic, injecting packets only between existing ones leaves the
total gap fixed and raises the number of gaps, so the mean gap can only fall; allowing injection after the last packet
lengthens the total gap, and the mean gap moves both ways.

\subsection{Robust Detectors}
\label{sec:fw:robustness}

Gimenez et al.\ prove that a detector is robust when the attacker's edits can only raise each feature and the score
never falls when features rise~\cite{gimenez2025certifiably}; a feature that can only fall is covered by negating it. With a computed table, this becomes a
corollary of Theorem~\ref{thm:fw:directions}.

\corollary[Robustness]
\label{cor:fw:robustness}
{
Let $D$ be sound for $\mathcal{M}$ over a closed population $F$, and let the score $h$ read no feature marked both and
never fall when up features rise, down features fall, and fixed features stay. Then
$h(\varphi(z')) \ge h(\varphi(z))$ for every $z \in F$ and every $z' \in \mathrm{Reach}_\mathcal{M}(z)$.
}

\noindent A gradient-boosted tree model with monotone constraints is such a score. What Theorem~\ref{thm:fw:directions}
adds to Gimenez et al.'s route is the table, computed for every feature. Global program analyses, learned
representations, and time-decayed statistics fall outside the class (\Cref{sec:discussion}).

Several pieces of this section are classical. Proposition~\ref{prop:fw:side-effects}'s sums follow the delta rules
of incremental view maintenance~\cite{gupta1995maintenance,koch2014dbtoaster}, including the failure of the maximum
to have an inverse; the item table is a provenance annotation in the sense of Green et al.~\cite{green2007provenance};
recomputing only what an edit touches is the idea of self-adjusting computation~\cite{acar2002adaptive}; and the
directions are one step of the sign domain. What these lack, and what we add, is the bridge from an edit of the raw
input to a change of items: the locating positions and margins that say which items an edit of elements changes
without parsing again, the vocabularies of terms and readouts that a feature extractor of this shape is written in,
and signatures derived per input, which turn one sign computation into a table that holds over a population.
